\documentclass[a4paper]{article}
\usepackage{amsfonts}
\usepackage{amsmath}
\usepackage{amssymb}
\usepackage{graphicx}     

\begin{document}
  
\noindent \large Auteur: Niels Hertoghs \\
\noindent \large Studierichting: Informatica\\
\noindent \large Oefeningenreeks 2D nr. 10.14\\

\medskip

\normalsize

$\lim\limits_{x \rightarrow 2}  \dfrac{x^3-x^2-8x+12}{3x^3-10x^2+4x+8} = \dfrac{0}{0} $\\\\


\begin{tabular}{c|ccc|c}
			& $1$	& $-1$ 	& $-8$ 	& 			$12$ \\
	   $-3$ & 		& $-3$ 	& $12$ 	& 			$-12$ \\ \hline
			& $1$ 	& $-4$ 	& $4$ 	& 	$0$  \\

\end{tabular}\\

$\Rightarrow (x^2-4x+4) = (x-2)^2$\\

$=\lim\limits_{x \rightarrow 2}  \dfrac{(x+3)(x-2)^2}{3x^3-6x^2-4x^2+8x-4x+8} = \lim\limits_{x \rightarrow 2}  \dfrac{(x+3)(x-2)^2}{3x^2(x-2)-4x(x-2)-4(x-2)}$\\

$=\lim\limits_{x \rightarrow 2}  \dfrac{(x+3)(x-2)}{3x^2-4x-4}$\\\\

\begin{tabular}{c|cc|c}
	 & $3$	& $-4$ 	&		$-4$\\
$2$ & 		& $6$ 	& 		$4$ \\ \hline
 	 & $3$ 	& $2$ 	& $0$ \\
	
\end{tabular}\\

$=\lim\limits_{x \rightarrow 2}  \dfrac{(x+3)(x-2)}{(3x+2)(x-2)} = \lim\limits_{x \rightarrow 2}\dfrac{(x+3)}{(3x+2)} = \dfrac{5}{8}$\\


\begin{thebibliography}{999}
%\bibitem{a} Referentie
\end{thebibliography}
\end{document} 
